\chapter{Methodology} \label{Chap3}\label{chap:methodology}

\section{Research scope and methodological rationale}

The empirical design evaluates six linked components: London data preparation,
large-image density and colour inference, independent-mesh processing, static
and on-demand ChordAtlas integration, a matched representation comparison, and
component verification.  Together,
they test whether overlapping satellite-conditioned predictions can form a
continuous coloured surface and be converted into coordinate-explicit,
building-level assets for an interactive procedural environment.

A contiguous Sat3DGen experiment measures overlap fusion and global colour
recovery.  A separate experiment uses retained OBJ outputs to isolate
geospatial placement, mesh stitching, OSM/DSM processing, and building
extraction.  Both routes are evaluated at the ChordAtlas publication boundary.
Data-quality tests characterise the London multimodal collection.  Retained
ChordAtlas appearance exports provide representation-level evidence.

\section{Systems evaluation design}

The evaluation combines data-integrity analysis, coordinate tests, staged
geometry measurements, controlled regression cases, topological inspection,
and cross-language integration tests.  Source tracing establishes the
transformations applied by the active pipeline; generated artefacts and
manifests provide stage-level measurements; and executable checks test the
coordinate, geometry, cache, and publication contracts.  The ChordAtlas bridge
is compared with the existing workflow at the level of representation,
selection, conversion, and publication behaviour.

The research was organised around five evaluation questions:

\begin{description}
  \item[EQ1] Are the source data and pairing conventions internally consistent
  enough to support the intended cross-view and geospatial uses?
  \item[EQ2] How do density-coordinate placement, overlap weighting, and
  second-pass colour recovery affect boundary continuity, colour completeness,
  and resource use in large-image inference?
  \item[EQ3] How did the executed mesh-processing stages affect mesh structure, continuity, and height, and was their observed behaviour consistent with the intended design?
  \item[EQ4] What do the static and on-demand paths demonstrably achieve in
  ChordAtlas, and where do their atomicity and cache boundaries lie?
  \item[EQ5] How do the retained CityEngine and ChordAtlas artefacts differ in
  representation from the two matched Sat3DGen routes, and which comparisons
  remain valid without geometric or appearance ground truth?
\end{description}

\section{Evaluated configuration}

The evaluated configuration comprises the London data-preparation workflow,
the released Sat3DGen checkpoint and large-image extension, the active
independent-mesh package, and the ChordAtlas integration layer.  Runtime imports
and entry-point tracing identify the components used by each experiment.
Functional analysis covers GUI entry points, worker and process boundaries,
stable selection identities, manifest fields, cache validation, publication
order, rollback, and validation depth.  Reproduction details are recorded in
Appendix~\ref{app:evidence-record}.

\section{System overview}

Figure~\ref{fig:system-overview} presents the conceptual data flow derived from
the active implementation.  Chapter~\ref{chap:results} assigns execution and
validation status to the individual paths.

\begin{figure}[H]
\centering
\resizebox{0.98\textwidth}{!}{%
\begin{tikzpicture}[
  x=1mm,
  y=1mm,
  box/.style={draw=uclblue,rounded corners,align=center,minimum width=27mm,
    minimum height=11mm,text width=24mm,fill=lightgrey},
  data/.style={draw=evidencegreen,rounded corners,align=center,minimum width=25mm,
    minimum height=10mm,text width=22mm,fill=evidencegreen!8},
  context/.style={draw=evidencegreen,rounded corners,align=center,minimum width=30mm,
    minimum height=10mm,text width=27mm,fill=evidencegreen!8},
  arrow/.style={-{Latex[length=2.3mm]},thick,uclblue}
]
\node[data] (builder) at (0,22) {London imagery\\and metadata};
\node[box] (sat) at (31,22) {Overlapping Sat3DGen\\window inference};
\node[box] (density) at (63,22) {Fractional density\\fusion and one MC};
\node[box] (colour) at (95,22) {Second-pass\\vertex RGB fusion};

\node[data] (cached) at (0,0) {Cached Sat3DGen\\OBJ tiles};
\node[box] (place) at (31,0) {Coordinate placement\\and crop};
\node[box] (stitch) at (63,0) {Stitching and\\mesh cleanup};
\node[box] (publish) at (95,0) {Building extraction\\and publication};
\node[box] (chord) at (127,0) {ChordAtlas\\renderers and MiniMesh};
\node[context] (geo) at (95,-22) {OSM footprints and\\EPSG:27700 DSM};

\draw[arrow] (builder) -- (sat);
\draw[arrow] (sat) -- (density);
\draw[arrow] (density) -- (colour);
\draw[arrow] (colour) -- (publish);
\draw[arrow] (cached) -- (place);
\draw[arrow] (place) -- (stitch);
\draw[arrow] (stitch) -- (publish);
\draw[arrow] (geo) -- (publish);
\draw[arrow] (publish) -- (chord);
\end{tikzpicture}%
}
\caption[Conceptual architecture of the two evaluated processing routes.]{Conceptual project architecture.  The upper route performs contiguous
density-field fusion and colour recovery; the lower route evaluates
independent-mesh assembly.  Both converge on building extraction and
publication before ChordAtlas.  OSM footprints and the DSM enter from below as
independent geospatial inputs rather than sequential processing stages.}
\label{fig:system-overview}
\end{figure}

\section{London data preparation}

The London data builder is evaluated as a staged data contract supporting
cross-view pairing and downstream geospatial processing.  The protocol measures
coverage, duplication, partition leakage, label geometry, and auxiliary-channel
consistency.
Figure~\ref{fig:data-builder-flow} links each processing stage to its available
artefacts and audit denominator.

\begin{figure}[H]
\centering
\resizebox{0.98\textwidth}{!}{%
\begin{tikzpicture}[
  node distance=8mm and 7mm,
  source/.style={draw=evidenceblue,rounded corners,align=center,
    minimum height=11mm,text width=29mm,fill=evidenceblue!8},
  stage/.style={draw=uclblue,rounded corners,align=center,
    minimum height=11mm,text width=29mm,fill=lightgrey},
  warn/.style={draw=evidenceorange,rounded corners,align=center,
    minimum height=11mm,text width=31mm,fill=evidenceorange!8},
  arrow/.style={-{Latex[length=2mm]},thick,uclblue}
]
\node[source] (points) {OSM-highway\\query points};
\node[stage,right=of points] (match) {Panorama\\matching};
\node[stage,right=of match] (legacy) {Latitude-based\\data splits};
\node[stage,right=of legacy] (grid) {Web Mercator\\tile grid};
\node[stage,right=of grid] (overlap) {Overlapping\\grid variants};
\draw[arrow] (points) -- (match);
\draw[arrow] (match) -- (legacy);
\draw[arrow] (match) -- (grid);
\draw[arrow] (grid) -- (overlap);

\node[source,below=15mm of match] (pano) {Panorama branches\\legacy + equirectangular};
\node[warn,right=of pano] (labels) {VIGOR-like labels\\4 candidates/query};
\node[warn,right=of labels] (channels) {Auxiliary channels\\DSM / OSM\\sky / relative depth};
\node[source,right=of channels] (layouts) {London and Sat3DGen\\directory layouts};
\draw[arrow] (match) -- (pano);
\draw[arrow] (grid) -- (labels);
\draw[arrow] (grid) -- (channels);
\draw[arrow] (overlap) |- (layouts);
\draw[arrow] (pano) -- (labels);
\draw[arrow] (labels) -- (channels);
\draw[arrow] (channels) -- (layouts);
\end{tikzpicture}%
}
\caption{Data-builder stages and the contracts tested for completeness and
geometric consistency.}
\label{fig:data-builder-flow}
\end{figure}

\subsection{Point matching, tiling, and imagery}

The data-builder begins from WGS84 point features derived from an OSM highway
query.  A manually defined longitude--latitude region and a zoom-20 Web
Mercator grid assign satellite centres.  A panorama is matched when its point
lies inside a satellite tile's half-width and half-height.  Chapter~\ref{chap:results}
reports the resulting population and coverage.

For longitude $\lambda$ in degrees, latitude $\phi$ in radians, zoom $z$, and
world-pixel scale $S=256\,2^z$, the scripts use the standard Web Mercator form

\begin{align}
p_x &= S\frac{\lambda+180}{360},\\
p_y &= S\left[\frac{1}{2}-\frac{1}{4\pi}
  \ln\left(\frac{1+\sin\phi}{1-\sin\phi}\right)\right].
\end{align}

Latitude is clamped to the Web Mercator limit.  The nominal ground resolution
at latitude $\phi$ is

\begin{equation}
r(\phi,z)=\frac{2\pi R\cos\phi}{256\,2^z},
\end{equation}

where $R$ is the Web Mercator sphere radius.  At $\phi=51.5074^\circ$ and
$z=20$, the implementation gives approximately
\SI{0.093}{\metre\per\pixel}, so a 640-pixel tile spans about
\SI{59.5}{\metre}.  Values used in the audit were recalculated from the
coordinate metadata.

The acquisition procedure requests 640-by-320 Street View images with a
120-degree field of view and 640-by-640 satellite images, nominally at zoom 20.
The legacy branch does not retain the returned panorama identifier or snapped
coordinates.  The equirectangular branch requests six cube faces and projects them into 1,536-by-768
equirectangular panoramas.  The projection and image dimensions were treated
as data transformations, while observation independence was tested separately
through content hashes.

\subsection{Splits and candidate labels}

The legacy training, validation, and test split sorts by latitude and assigns
70, 15, and 15 per cent respectively.  Spatial sorting can reduce arbitrary
mixing, but there is no boundary buffer, and duplicate returned panoramas can
cross a boundary.

The VIGOR-style branch fills a 16-by-17 Cartesian grid of observed
latitude and longitude centres and chooses the four nearest satellite centres
for every panorama.  For candidate $j$, an offset was classified as containing
the query only when both absolute pixel components were at most 320.  The
audit recorded the count outside this square by candidate rank and the maximum
absolute component.  This tests geometric compatibility rather than mere
four-column syntactic compatibility with VIGOR \cite{zhu2021vigor}.

Two denser satellite grids were also reconstructed.  A 320-pixel stride gives
approximately 50 per cent overlap, while a 576-pixel stride gives 10 per cent
overlap.  Coverage calculations account for coordinate rounding and incomplete
placement at non-divisible outer boundaries.

\subsection{Auxiliary channels}

The pipeline includes binary sky masks, model-derived relative depth, a
DSM-derived normalised raster, and OSM semantic layers.  The DSM branch
transforms WGS84 coordinates into the raster CRS, reprojects a patch with
bilinear sampling, and scales its 5th--95th percentile range independently.
Consequently, it was evaluated as a normalised surface-height channel, not
metric depth.

For OSM, raw element types and tags were counted separately from raster output.
Supporting nodes returned by recursive Overpass queries were identified by
their lack of category tags.  Multipolygon reconstruction and inner-ring
assignment were traced through geometry construction and drawing.  This made
it possible to distinguish an existing image from a semantically valid mask.

\section{Coordinate and vertical-frame contract}

The bridge defines a \gls{localframe} at $(\lambda_0,\phi_0)$ using
$M=111{,}320$ metres per latitude degree:

\begin{align}
x &= (\lambda-\lambda_0)M\cos\phi_0,\\
z &= -(\phi-\phi_0)M.
\end{align}

Thus X points east, Z points south, and Y is vertical.  The inverse uses the
same constants, so algebraic round trips should be exact to floating-point
precision.  The frame is illustrated in Figure~\ref{fig:coordinate-frame}.

\begin{figure}[H]
\centering
\begin{tikzpicture}[scale=1.0,>=Latex]
  \coordinate (o) at (0,0);
  \draw[->,very thick,evidenceblue] (o) -- (4.2,0) node[right]{X: east};
  \draw[->,very thick,evidencegreen] (o) -- (0,3.0) node[above]{Y: height};
  \draw[->,very thick,evidenceorange] (o) -- (-2.2,-1.5) node[below left]{Z: south};
  \fill (o) circle (2.2pt) node[above right]{stored WGS84 origin};
  \draw[dashed,gray] (-3.0,0) -- (4.8,0);
  \draw[dashed,gray] (0,-2.0) -- (0,3.5);
  \node[align=left,anchor=west] at (4.8,2.0)
    {$x=(\lambda-\lambda_0)M\cos\phi_0$\\
     $z=-(\phi-\phi_0)M$};
\end{tikzpicture}
\caption{Local coordinate contract used by the bridge.  The reversible local
approximation is convenient for block-scale processing but is not identical to
EPSG:27700.}
\label{fig:coordinate-frame}
\end{figure}

Coordinate evaluation had two parts.  First, the AOI centre and four corners
were transformed to local coordinates and back, measuring residual in metres.
Secondly, the same points were transformed from EPSG:4326 to EPSG:27700 using
\code{pyproj}; projected easting/northing differences from the origin were
compared with local X/south values.  This separates inverse consistency from
external metric agreement.

The vertical contract has three distinct quantities: generated-mesh local Y,
DSM surface elevation, and a workspace vertical offset.  Their values were
read from manifests and paired OBJ vertices.

\section{Large-image and mesh-processing protocols}

\subsection{Rectangular density-field fusion}

The large-image extension retains the released Sat3DGen checkpoint and its
standard 256-by-256 network transform.  It accepts a rectangular source mosaic
without first truncating it to a central square.  In the evaluated setting, a
640-pixel source window moves by 160 pixels, giving 75 per cent nominal overlap;
each source window is resized by the unchanged model transform.  The
\code{mesh resolution} parameter is the volumetric sampling size, not the image
input size.  A value of $N=192$ and the checkpoint's 0.8 position scale produce
a cropped density patch of $154\times154\times173$ samples with a 19-voxel
border.

Let $s_i$ be the source-image start of window $i$ on either horizontal axis and
$L=154$ the cropped density width.  Its continuous density origin is

\begin{equation}
o_i=s_i\frac{L}{640}.
\end{equation}

The origin is retained as a floating-point value rather than rounded to an
integer cell.  Each density sample is bilinearly distributed to the four
neighbouring cells of the global horizontal grid.  A separable raised-cosine
window tapers the interior 19-voxel boundary of each patch; an edge coincident
with the outer mosaic boundary retains unit weight.  For a global density
position $p$, the accumulated field is

\begin{equation}
\bar{D}(p)=\frac{\sum_i w_i(p)D_i(p)}{\sum_i w_i(p)},
\label{eq:density-feather-fusion}
\end{equation}

where $D_i$ is the resampled prediction and $w_i$ combines its bilinear
placement and spatial taper.  A zero denominator is treated as an invalid
coverage state.  Marching Cubes is applied once to $\bar{D}$ at the evaluated
isovalue of 4.5, so overlap is resolved before surface extraction.

Algorithm~\ref{alg:fractional-density-fusion} gives the corresponding
semi-formal procedure.  The two accumulators separate predicted density from
its support weight, which permits a final normalised field to be extracted by
a single Marching Cubes operation.

\begin{algorithm}[H]
\caption{Fractional feather fusion of overlapping density predictions}
\label{alg:fractional-density-fusion}
\begin{algorithmic}[1]
\Require Rectangular mosaic $I$; source-window starts $S$; predictor $M$;
cropped density width $L$; halo $h$; isovalue $\tau$
\Ensure Global vertices $V$ and faces $F$
\State Initialise density sum $A\gets 0$ and horizontal weight sum $W\gets 0$
\ForAll{$(r_i,c_i)\in S$}
  \State $D_i\gets\Call{PredictDensity}{M,I,r_i,c_i}$
  \State $D_i\gets\Call{CropHalo}{D_i,h}$
  \State $(o_{r},o_{c})\gets(r_iL/640,c_iL/640)$
  \State $K_i\gets\Call{RaisedCosineWeights}{L,h,\text{mosaic-edge flags}}$
  \State $(f_r,f_c)\gets(o_r-\lfloor o_r\rfloor,o_c-\lfloor o_c\rfloor)$
  \ForAll{$(a,b)\in\{0,1\}^{2}$}
    \State $\alpha_{ab}\gets(af_r+(1-a)(1-f_r))(bf_c+(1-b)(1-f_c))$
    \State \Call{Accumulate}{$A,\alpha_{ab}K_iD_i,\lfloor o_r\rfloor+a,
      \lfloor o_c\rfloor+b$}
    \State \Call{Accumulate}{$W,\alpha_{ab}K_i,\lfloor o_r\rfloor+a,
      \lfloor o_c\rfloor+b$}
  \EndFor
\EndFor
\State \textbf{require} $W(p)>0$ for every covered horizontal position $p$
\State $\bar D\gets A/W$ \Comment{broadcast $W$ over the vertical axis}
\State $(V,F)\gets\Call{MarchingCubes}{\bar D,\tau}$
\State \Return $(V,F)$
\end{algorithmic}
\end{algorithm}

\subsection{Memory-bounded vertex-colour recovery}

The released large-image wrapper discards each tri-plane after density
prediction and therefore exports geometry without the vertex-colour query used
by single-image inference \cite{qianSat3DGenCode2026}.  The added colour pass
keeps the extracted global vertices and faces fixed.  It reprocesses one source
window at a time, reconstructs that window's tri-plane, identifies the global
vertices supported by the same fractional feather weights, and queries the
pretrained colour network in batches.

The global Marching-Cubes axes are $[\mathrm{row},\mathrm{column},z]$.  For a
window with continuous origin $(o_r,o_c)$, a supported vertex is mapped back to
the full $192^3$ model grid as

\begin{equation}
q=\left[\mathrm{column}-o_c+19,
         \mathrm{row}-o_r+19,
         z+19\right]\frac{2}{191}-1.
\label{eq:large-image-colour-coordinate}
\end{equation}

The horizontal query is clamped only for the fractional halo outside the
cropped density support; the support weight itself is evaluated without that
clamp.  With $c_i(v)$ denoting the RGB query for vertex $v$, colour is blended
in floating point as

\begin{equation}
\bar{C}(v)=\frac{\sum_i w_i(v)c_i(v)}{\sum_i w_i(v)}.
\label{eq:colour-feather-fusion}
\end{equation}

A two-dimensional spatial index limits each window to nearby vertices, and
only one tri-plane is resident at a time.  Export requires positive colour
weight for every vertex, finite RGB values, unchanged vertex and face arrays,
and an exact colour round trip after reloading the PLY.

Algorithm~\ref{alg:vertex-colour-recovery} describes the second pass.  It uses
the geometry-fusion weights again, but accumulates RGB only for vertices whose
continuous coordinates lie inside a window's weighted support.

\begin{algorithm}[H]
\caption{Memory-bounded recovery of colour on a fixed global mesh}
\label{alg:vertex-colour-recovery}
\begin{algorithmic}[1]
\Require Mosaic $I$; fixed mesh $(V,F)$; window plan $S$; predictor $M$;
fusion metadata
\Ensure Per-vertex RGB array $C$ on unchanged $(V,F)$
\State Build a two-dimensional spatial index $B$ over the horizontal axes of $V$
\State Initialise RGB sum $R\gets 0$ and vertex weight sum $Z\gets 0$
\ForAll{windows $i\in S$}
  \State $T_i\gets\Call{EncodeWindow}{M,I,i}$ \Comment{one resident tri-plane}
  \State $(J,w)\gets\Call{SupportedVertices}{V,B,i,\text{fusion metadata}}$
  \ForAll{paired batches $(Q,w_Q)$ drawn from $(J,w)$}
    \State $q\gets\Call{ModelCoordinates}{V[Q],i}$
    \State $c\gets\Call{QueryColour}{M,T_i,q}$
    \State $R[Q]\gets R[Q]+w_Qc$
    \State $Z[Q]\gets Z[Q]+w_Q$
  \EndFor
  \State Release $T_i$
\EndFor
\State \textbf{require} $Z(v)>0$ and finite for every $v\in V$
\State $C\gets\Call{QuantiseRGB}{R/Z}$
\State \Call{ExportAndReloadVerify}{$V,F,C$}
\State \Return $C$
\end{algorithmic}
\end{algorithm}

\begin{figure}[H]
\centering
\resizebox{0.98\textwidth}{!}{%
\begin{tikzpicture}[
  node distance=7mm and 6mm,
  source/.style={draw=evidenceblue,rounded corners,align=center,
    text width=25mm,minimum height=10mm,fill=evidenceblue!8},
  stage/.style={draw=uclblue,rounded corners,align=center,
    text width=27mm,minimum height=10mm,fill=lightgrey},
  output/.style={draw=evidencegreen,rounded corners,align=center,
    text width=27mm,minimum height=10mm,fill=evidencegreen!8},
  arrow/.style={-{Latex[length=2mm]},thick,uclblue}
]
\node[source] (mosaic) {Rectangular zoom-20 mosaic};
\node[stage,right=of mosaic] (windows) {Overlapping source windows};
\node[stage,right=of windows] (patches) {Cropped density patches};
\node[stage,right=of patches] (fusion) {Fractional splat and feather};
\node[output,right=of fusion] (field) {Global density and one MC};
\draw[arrow] (mosaic) -- (windows);
\draw[arrow] (windows) -- (patches);
\draw[arrow] (patches) -- (fusion);
\draw[arrow] (fusion) -- (field);

\node[stage,below=11mm of windows] (triplane) {Recomputed tri-plane per window};
\node[stage,right=of triplane] (query) {Supported-vertex colour query};
\node[stage,right=of query] (blend) {Weighted floating-point RGB};
\node[output,right=of blend] (mesh) {Fixed geometry with vertex colour};
\draw[arrow] (windows) -- (triplane);
\draw[arrow] (triplane) -- (query);
\draw[arrow] (query) -- (blend);
\draw[arrow] (field) -- (mesh);
\draw[arrow] (blend) -- (mesh);
\end{tikzpicture}%
}
\caption{Two-pass large-image method.  The first pass blends density before a
single surface extraction; the second revisits each window to blend colour on
the unchanged global vertices.}
\label{fig:large-image-method-flow}
\end{figure}

\subsection{Independent-mesh stages and relation to Sat3DGen}

The evaluated sequence loads configuration, OSM, and DSM data; discovers OBJ
tiles; parses, crops, transforms, classifies, and pre-aligns each tile; removes
lower geometry; stitches vertex groups; applies DSM correction; and exports
paired DSM and no-DSM building branches.  The experiment isolates these
post-inference operations from upstream generation and optional fa\c{c}ade
processing.

Table~\ref{tab:sat3dgen-method-comparison} distinguishes the released
large-image wrapper, the density/colour extension, and the independent-mesh
geospatial path.

\begin{table}[H]
\centering
\caption{Relationship among the released Sat3DGen large-image wrapper, the
density/colour extension, and mesh-space geospatial integration.}
\label{tab:sat3dgen-method-comparison}
\footnotesize
\begin{tabularx}{\textwidth}{p{0.16\textwidth}XXX}
\toprule
Aspect & Released large-image wrapper & Density/colour extension &
Mesh-space and ChordAtlas path\\
\midrule
Learned representation & Released DINOv3-conditioned tri-plane density field &
Same checkpoint and model transform; no retraining & Released model or cached
OBJ enters explicit local coordinates\\
Overlap & Integer-aligned cropped fields and equal-weight averaging &
Fractional origins, bilinear placement, and raised-cosine weights & Independent
OBJ tiles are cropped, placed, and stitched\\
Surface and colour & One global Marching Cubes surface; large-image PLY has no
colour & One global surface followed by weighted per-vertex RGB recovery & RGB
OBJ assets, building extraction, and a separate ChordAtlas colour renderer\\
Geospatial context & Normalised scene coordinates & Rectangular mosaic and
explicit image-to-density scale & X-east/Y-up/Z-south frame, OSM footprints,
and EPSG:27700 DSM\\
\bottomrule
\end{tabularx}
\end{table}

Each stage was quantified using OBJ header counts and, where necessary,
streamed vertex comparison.  The following metrics were computed:

\begin{itemize}
  \item vertices and faces before crop, after crop, after bottom removal, after
  stitching, and after building extraction;
  \item geographic cluster count and separation from filename coordinates;
  \item number and distribution of per-vertex DSM $\Delta Y$ values;
  \item finite vertices, valid indices, zero-area and duplicate faces, boundary
  and non-manifold edges, connected components, and edge-incidence
  watertightness for a representative paired building; 
\end{itemize}

\subsection{Controlled tests of geometry rules}

Two small synthetic fixtures tested whether documentation and configuration
were enforced.  For bottom removal, one triangle had only one vertex within
\SI{0.5}{\metre} of the global minimum and a second triangle was fully elevated.
The stated all-three-vertices rule predicts that neither should be removed.  The
active function was invoked without modifying source code and its returned
face set recorded.

For stitching, two identically wound triangles had coincident X/Z coordinates
but were separated by \SI{10}{\metre} in Y.  They were assigned different tile
ranges and stitched with an X/Z radius of \SI{0.01}{\metre}.  The experiment
recorded vertex/face counts and retained attributes.  It directly tests whether
the configured maximum Y difference participates in matching.

\section{ChordAtlas integration protocols}

\subsection{Static workspace}

The static bridge combines a target bounding box, geospatial footprints,
existing mesh tiles, panoramas, and DSM prerequisites in a staged workspace.
Python records the manifest and local-frame footprint geometry; Java constructs
the ChordAtlas object graph and converts OBJ geometry into indexed MiniMesh
tiles.  A stored source offset preserves the original centroid during
conversion instead of implicitly recentering the input.

Validation checks required files, footprint geometry, MiniMesh structure, and
coordinate metadata.  Mesh coverage was independently intersected with the
target AOI. 

\subsection{On-demand selection and publication}

The on-demand path is embedded in the existing ChordAtlas GIS generator.  A
selected polygon is canonicalised by
ring start, winding, holes, duplicate removal, and six-decimal rounding before
generating a stable selection identifier.  Processing is serialised on one
worker, and interface updates return to the GUI event thread.

The Java service writes an atomic request file and launches the Python bridge.
The request distinguishes independent satellite-tile processing from the
large-image mode.  The tile path derives a geographically anchored allow-list
and selects fresh or cached corrected geometry.  The large-image path plans a
rectangular zoom-20 mosaic, executes or validates the density and colour
stages, and converts the fused mesh
once into the project local frame.  Both routes apply DSM processing and
associate geometry with requested footprints.

Each publishable building receives an exact source-mesh manifest, geometry
after ground processing, and metadata in a staging tree.  Geometry ownership
and filtering are applied per footprint rather than reusing one
selection-wide mesh under several building identities.  In large-image mode,
Python additionally requires RGB on every OBJ vertex and verifies that DSM
processing preserves colour indices.  Java displays this OBJ through a
dedicated vertex-colour renderer; the per-building MiniMesh is generated as a
separate representation used by the ChordAtlas workspace.  The staged Python
publication is installed atomically within its own boundary and marked
complete before Java conversion and renderer construction.

Figure~\ref{fig:publication-lifecycle} makes the transaction boundary explicit.

\begin{figure}[H]
\centering
\begin{tikzpicture}[
  node distance=7mm,
  state/.style={draw=uclblue,rounded corners,align=center,text width=38mm,
    minimum height=9mm,fill=lightgrey},
  good/.style={draw=evidencegreen,rounded corners,align=center,text width=38mm,
    minimum height=9mm,fill=evidencegreen!8},
  arrow/.style={-{Latex[length=2mm]},thick}
]
\node[state] (select) {GIS polygon and stable ID};
\node[state,right=of select] (python) {Python plan, inputs, generation, DSM, extraction};
\node[good,right=of python] (ready) {Python publication\\complete};
\node[state,below=of ready] (mini) {Java per-building MiniMesh conversion};
\node[good,left=of mini] (gui) {GUI building-generator creation};
\draw[arrow] (select) -- (python);
\draw[arrow] (python) -- (ready);
\draw[arrow] (ready) -- node[right,font=\scriptsize]{outside Python commit} (mini);
\draw[arrow] (mini) -- (gui);
\draw[dashed,evidenceorange,thick] ($(ready.south west)+(-1mm,-3mm)$) --
  ($(ready.south east)+(1mm,-3mm)$) node[midway,below=1mm,font=\scriptsize]
  {atomicity boundary};
\end{tikzpicture}
\caption{On-demand publication lifecycle.  Python completion precedes
the Java conversion required by ChordAtlas, so end-to-end publication is not a
single atomic transaction.}
\label{fig:publication-lifecycle}
\end{figure}

\subsection{Selection-scoped appearance inputs}

The appearance branch reuses the same geographic selection and local frame as
the building publication.  For roofs, the bridge selects the satellite tiles
already bound to the request, mosaics them in north-up orientation, clips them
to the building footprint, and publishes a reference image and style
mask.  Existing ready references can be added to an earlier building
publication without modifying its geometry.  ChordAtlas validates the
reference metadata before exposing it to the fa\c{c}ade and roof style
interface.

For street-level context, candidate viewpoints are sampled around the selected
OSM footprint and converted to WGS84.  Street View metadata are resolved and
deduplicated before acquisition.  The service presents one sample for approval,
downloads the remaining accepted batch, promotes the complete set, and writes
the panorama positions into ChordAtlas's local coordinate frame.  Figure~\ref{fig:appearance-input-flow}
summarises the two paths.

\begin{figure}[H]
\centering
\resizebox{0.98\textwidth}{!}{%
\begin{tikzpicture}[
  node distance=7mm and 6mm,
  source/.style={draw=evidenceblue,rounded corners,align=center,
    text width=27mm,minimum height=10mm,fill=evidenceblue!8},
  stage/.style={draw=uclblue,rounded corners,align=center,
    text width=29mm,minimum height=10mm,fill=lightgrey},
  output/.style={draw=evidencegreen,rounded corners,align=center,
    text width=30mm,minimum height=10mm,fill=evidencegreen!8},
  arrow/.style={-{Latex[length=2mm]},thick,uclblue}
]
\node[source] (sat) {Selection-bound\\satellite tiles};
\node[stage,right=of sat] (mosaic) {North-up mosaic\\and footprint mask};
\node[output,right=of mosaic] (roof) {Hashed roof\\reference};
\node[stage,right=of roof] (style) {ChordAtlas style\\input contract};
\draw[arrow] (sat) -- (mosaic);
\draw[arrow] (mosaic) -- (roof);
\draw[arrow] (roof) -- (style);

\node[source,below=11mm of sat] (footprint) {Selected OSM\\footprint};
\node[stage,right=of footprint] (viewpoints) {Viewpoint sampling\\and metadata deduplication};
\node[stage,right=of viewpoints] (approval) {Sample approval\\and batch acquisition};
\node[output,right=of approval] (panos) {Local-frame\\panorama index};
\draw[arrow] (footprint) -- (viewpoints);
\draw[arrow] (viewpoints) -- (approval);
\draw[arrow] (approval) -- (panos);
\end{tikzpicture}%
}
\caption{Selection-scoped appearance-input workflows.  The roof path prepares
a footprint-matched satellite reference for the style interface; the panorama
path acquires and registers street-level observations in the same local frame.}
\label{fig:appearance-input-flow}
\end{figure}

Selection logs and result manifests were classified by processing state:
request rejection, planning, acquisition, inference, corrected-scene reuse,
building publication, Java conversion, and GUI creation. 

\section{Matched representation comparison}

The directly matched comparison fixes OSM way~369245408 as the main-building
identity.  The large-image and independent-tile requests contain the same
target bounds and the same two footprint geometries.  The former begins from
one fused density field; the latter begins from nine independently extracted
meshes and adds OSM pre-alignment, lower-surface removal, stitching, and DSM
processing.  Building-level vertices, faces, relief, edge incidence, and
vertex-colour variation are compared under this shared spatial request.

Cross-system inspection is restricted to representation.  A retained
CityEngine 2026 export records one named initial shape, UV coordinates, and
UV-mapped materials.  The two
route-labelled ChordAtlas exports have envelopes consistent with the main
footprint; the standalone
ChordAtlas example has matched location.  All examples are
rendered with a common orthographic camera and metric scale within each figure.
The audit records polygon or triangle representation, vertex colour, texture
coordinates, and diffuse/normal/specular map organisation.  These properties
are not treated as fidelity scores because no surveyed geometry or registered
street-level appearance reference is available.

\section{Verification protocol}

\subsection{Executable tests}

Component tests cover configuration, coordinate frames, footprints, DSM
prerequisites, OBJ inspection, tile planning, image integrity, manifests,
cache freshness, publication, rollback, MiniMesh conversion, OBJ slicing, and
vertex-colour parsing.  Separate large-image regression cases exercise
fractional placement, non-zero fusion coverage, seam continuity, colour-query
coordinates, weighted RGB, and PLY round trips.  The bridge suite also covers
large-image selection, roof-reference and panorama planning, promotion, and
coordinate transfer.  Python, algorithm, and Java tests are reported
separately because they exercise different runtime and data-transfer
boundaries.  Detailed results are provided in
Appendix~\ref{app:evidence-record}.

\subsection{Validity review}

Construct validity was assessed by matching each measure to the property it
was intended to represent.  The review distinguishes geographic extent from
triangle coverage, file availability from semantic correctness, topological
validity from geometric accuracy, and producer-side publication from
consumer-side conversion.  External validity is examined through the spatial
and dataset coverage of the case study.
